\documentclass[10pt]{article}
\usepackage[a4paper,margin=20mm]{geometry}
\usepackage[T1]{fontenc}
\usepackage{lmodern,microtype}
\usepackage{amsmath,amssymb,mathtools,bm}
\usepackage{booktabs,tabularx,array}
\usepackage[dvipsnames]{xcolor}
\usepackage[hidelinks]{hyperref}
\usepackage{enumitem}
\setlist{nosep,leftmargin=1.5em}
\setlength{\parindent}{0pt}
\setlength{\parskip}{5pt}
\definecolor{navy}{HTML}{183B56}
\definecolor{green}{HTML}{257A55}
\definecolor{soft}{HTML}{F2F7F4}
\newcommand{\R}{\mathbb{R}}
\newcommand{\C}{\mathbb{C}}
\newcommand{\wrap}{\operatorname{wrap}}
\newcommand{\LN}{\operatorname{LN}}
\newcommand{\MLP}{\operatorname{MLP}}
\newcommand{\mean}{\operatorname{mean}}
\newcommand{\code}[1]{\texttt{#1}}
\title{\vspace{-1.4em}\color{navy}\textbf{GridSFM as an AC Power-Flow Surrogate}\\
\large Exact formulation of the released backbone and PPC adaptation}
\author{Implementation note for \code{train\_valid\_test\_gridsfm.py}}
\date{August 2026}

\begin{document}
\maketitle
\vspace{-1.5em}
\begin{center}
\fcolorbox{green}{soft}{\parbox{0.93\linewidth}{
\textbf{Scope.} These equations describe the released
\code{gridsfm.GridTransformerBackbone} driven by the AC-PF branch-row adapter.
GridSFM was originally structured for AC-OPF and ships pretrained weights; here
its voltage heads are fine-tuned or retrained against Newton PF labels. Optional
known-operator postprocessing is excluded.
}}
\end{center}

\section{Target AC-PF Mapping}

For \(\underline V_i=v_i e^{\mathrm j\theta_i}\), define
\begin{equation}
 \underline{\bm S}^{\rm calc}=\underline{\bm V}\odot
 \overline{\bm Y\underline{\bm V}},\qquad
 \Delta\bm P=\bm P^{\rm set}-\Re\underline{\bm S}^{\rm calc},\quad
 \Delta\bm Q=\bm Q^{\rm set}-\Im\underline{\bm S}^{\rm calc}.
\end{equation}
The surrogate learns
\begin{equation}
 f_{\Theta}(\mathcal G,\underline{\bm S}^{\rm set},\bm V^0)
 \longmapsto (\hat{\bm v},\hat{\bm\theta})
 \approx(\bm v^\star,\bm\theta^\star).
\end{equation}
Active residuals are evaluated on non-slack buses and reactive residuals on PQ
buses, matching the fixed-bus-type Newton system.

\section{PPC-to-GridSFM Heterogeneous Graph}

Before preprocessing, the adapter creates node sets
\begin{equation}
 \mathcal V_b,\quad \widetilde{\mathcal V}_g,\quad
 \mathcal V_d,\quad \mathcal V_s,
\end{equation}
for buses, proxy generators, loads, and shunts. The released preprocessing then
promotes branches to nodes and introduces cycle nodes:
\begin{equation}
 \mathcal V=\mathcal V_b\cup\widetilde{\mathcal V}_g\cup\mathcal V_d\cup
 \mathcal V_s\cup\mathcal V_{\ell}\cup\mathcal V_{\tau}\cup\mathcal V_c.
\end{equation}
Here \(\mathcal V_{\ell}\) and \(\mathcal V_{\tau}\) represent AC-line and
transformer branch rows, while \(\mathcal V_c\) represents a fundamental cycle
basis. Native GridSFM expects physical generator rows; the PPC PF adapter instead
creates one proxy generator per PV/REF bus.

\subsection{Raw node features}

For a bus \(i\),
\begin{equation}
 \bm x_i^b=[V_{n,i}^{\rm kV},\,t_i,\,v_i^{\min},\,v_i^{\max}]\in\R^4,
\end{equation}
where \(t_i\in\{1,2,3,4\}\) encodes PQ, PV, REF, or isolated. The start
magnitude is retained separately as \(v_{i}^{\rm set}=v_i^0\).

For proxy generator \(\tilde g_i\), load \(d_i\), and shunt \(s_i\),
\begin{align}
 \bm x_{\tilde g_i}^g&=[m_{\rm base},P_g,P_g^{\min},P_g^{\max},Q_g,
 Q_g^{\min},Q_g^{\max},v_g^{\rm set},c_2,c_1,c_0]\in\R^{11},\\
 \bm x_{d_i}^d&=[P_{d,i},Q_{d,i}]\in\R^2,\\
 \bm x_{s_i}^s&=[B_i^{\rm sh},G_i^{\rm sh}]\in\R^2.
\end{align}
In the PF adapter, generator limits are synthesized around bus-level injections,
costs are zero, and \(v_g^{\rm set}=v_i^0\). Reactive injection is set to zero at
PV/REF buses before graph construction because it is not specified by fixed-type
PF. Load channels are formed from the negative net injection and clipped at zero.

\subsection{Line and transformer features}

The adapter divides live branches into AC lines and transformers. An AC-line row is
\begin{equation}
 \bm e_{ij}^{\ell}=[\theta^{\min},\theta^{\max},b_{fr},b_{to},r,x,
 S_A,S_B,S_C]\in\R^9,
\end{equation}
and a transformer row is
\begin{equation}
 \bm e_{ij}^{\tau}=[\theta^{\min},\theta^{\max},r,x,S_A,S_B,S_C,
 \tau,\phi,b_{fr},b_{to}]\in\R^{11}.
\end{equation}
The series impedance is recovered from the PPC admittance,
\begin{equation}
 r+\mathrm jx=1/y_{ij},
\end{equation}
and transformers are identified from metadata, non-unit tap, phase shift, and
optionally endpoint nominal-voltage mismatch.

\section{Branch Promotion and Cycle Representation}

Each branch row becomes a node whose initial feature is its line or transformer
attribute. For branch \(e=(i,j)\), signed endpoint edges are created:
\begin{equation}
 i\xleftrightarrow{+1}e,\qquad j\xleftrightarrow{-1}e.
\end{equation}
Let \(B_1\) denote this oriented bus--branch incidence. A fundamental cycle
basis is built from a spanning tree. A cycle \(c\) has features
\begin{equation}
 \bm x_c^c=[|c|,\sum_{e\in c}|x_e|,\sum_{e\in c}|r_e|,
 \min_{e\in c}S_e^{\max}]\in\R^4,
\end{equation}
and signed cycle--branch incidence \(B_2(c,e)\in\{-1,0,+1\}\). This creates the
sequence
\begin{equation}
 \text{bus (0-form)}\longleftrightarrow
 \text{branch (1-form)}\longleftrightarrow
 \text{cycle (2-form)}.
\end{equation}

\section{Static Physics and Positional Features}

\subsection{Laplacian/DC features}

Using branch susceptances \(b_{ij}\approx1/x_{ij}\), GridSFM builds
\begin{equation}
 L_B=\sum_{(i,j)}b_{ij}(\bm e_i-\bm e_j)(\bm e_i-\bm e_j)^\top.
\end{equation}
A reduced DC solve supplies
\begin{equation}
 L_{B,\setminus r}\bm\theta_{\rm DC}=\bm P_{\setminus r},
\end{equation}
with the slack balancing total injection. Bus DC features are
\begin{equation}
 \bm p_i^{\rm DC}=[\theta_i^{\rm DC},\,\operatorname{mean}_{j\in\mathcal N(i)}
 |\theta_i^{\rm DC}-\theta_j^{\rm DC}-\phi_{ij}|,\,P_i^{\rm raw}]\in\R^3.
\end{equation}
Each branch receives
\begin{equation}
 [\Delta\theta_{ij}^{\rm DC},P_{ij}^{\rm DC},
 \min(|P_{ij}^{\rm DC}|/S_{ij}^{\max},2),\operatorname{sign}P_{ij}^{\rm DC}]\in\R^4.
\end{equation}
Five landmark effective-resistance moments are added to each bus. A four-column
bus-type one-hot is appended, giving \(4+5+3+4=16\) bus columns before learned
Hodge encoding.

\subsection{Learned Hodge positional encoding}

For each active form \(q\in\{0,1,2\}\), let \(L_q\) be the signed Hodge
operator induced by \(B_1\) and \(B_2\):
\begin{equation}
 L_0=B_1B_1^\top,\qquad
 L_1=B_1^\top B_1+B_2^\top B_2,\qquad
 L_2=B_2B_2^\top.
\end{equation}
After projecting to \(K_H=8\) channels, GridSFM performs \(T_H=4\) learned
diffusion steps:
\begin{equation}
 \bm z_q^{(0)}=W_q\bm x_q,\qquad
 \bm z_q^{(t+1)}=\bm z_q^{(t)}-\bm\alpha_{q,t}\odot L_q\bm z_q^{(t)}.
\end{equation}
The concatenated trajectory is mixed and normalized:
\begin{equation}
 \bm p_q^{\rm Hodge}=\LN\!\left(\MLP_q
 [\bm z_q^{(0)}\Vert\cdots\Vert\bm z_q^{(T_H)}]\right).
\end{equation}

An eight-dimensional graph context also summarizes system size, load/capacity
ratios, voltage-band width, ratings, and admittance scales:
\begin{equation}
 \bm g_{\rm ctx}\in\R^8.
\end{equation}

\section{GridTransformer Blocks}

Each node type is independently normalized and projected to hidden dimension
\(d=128\). The released default has \(L=8\) blocks and \(H=4\) heads.

\subsection{Per-type linear self-attention}

For each node type \(t\), attention is global within each graph but only among
nodes of that same type. With feature map \(\varphi(x)=\operatorname{ELU}(x)+1\),
\begin{align}
 Q&=W_Qh,\quad K=W_Kh,\quad Z=W_Vh,\\
 \bm M_{\gamma,h}&=\sum_{j\in\gamma,t}\varphi(K_j^h)(Z_j^h)^\top,
 &\bm k_{\gamma,h}&=\sum_{j\in\gamma,t}\varphi(K_j^h),\\
 \operatorname{Attn}_{t,h}(i)&=
 \frac{\varphi(Q_i^h)^\top\bm M_{\gamma(i),h}}
 {\varphi(Q_i^h)^\top\bm k_{\gamma(i),h}+\epsilon}.
\end{align}
This Performer-style expression is linear in the number of nodes rather than
dense quadratic attention. The residual update is
\begin{equation}
 h_t\leftarrow h_t+W_O\mathop{\Vert}_{h=1}^{H}\operatorname{Attn}_{t,h}.
\end{equation}

\subsection{Typed local message passing and FFN}

For signed incidence relations, GridSFM uses
\begin{equation}
 m_{u\to v}=s_{uv}W_{\rm msg}^{(r)}h_u,\qquad
 \tilde h_v=W_{\rm self}^{(r)}h_v+\operatorname{mean}_{u\in\mathcal N_r(v)}m_{u\to v}.
\end{equation}
For generator/load/shunt association relations it uses relation-specific
GraphSAGE mean aggregation. Relation outputs with the same target type are
summed by HeteroConv:
\begin{equation}
 h_v\leftarrow h_v+\sum_{r:\operatorname{dst}(r)=t(v)}\tilde h_v^{(r)}.
\end{equation}
Each block closes with a pre-normalized residual feed-forward network:
\begin{equation}
 h_v\leftarrow h_v+W_2^{t(v)}\operatorname{GELU}
 (W_1^{t(v)}\LN(h_v)).
\end{equation}

\section{Fusion and Terminal Heads}

After all blocks, relation-aware mean aggregates are fused into bus embeddings:
\begin{equation}
 h_i^{\rm grid,b}=\LN\!\left(W_bh_i^b+W_{\ell}\bar h_i^{\ell}
 +W_{\tau}\bar h_i^{\tau}+W_g\bar h_i^g+W_d\bar h_i^d
 +W_s\bar h_i^s+W_c\bar h_i^c\right).
\end{equation}
Generator embeddings fuse their associated bus. A global embedding is obtained
from mean and maximum pools of buses, line branches, transformer branches, and
cycles.

GridSFM applies its heads once, after fusion. Generator predictions are bounded
around their intervals:
\begin{align}
 \hat P_g&=\operatorname{clip}_{[P_g^{\min},P_g^{\max}]}
 \left(\tfrac12(P_g^{\min}+P_g^{\max})+H_P(h_g)\right),\\
 \hat Q_g&=\operatorname{clip}_{[Q_g^{\min},Q_g^{\max}]}
 \left(\tfrac12(Q_g^{\min}+Q_g^{\max})+H_Q(h_g)\right).
\end{align}
\par\nopagebreak[4]
\textbf{Voltage-angle head.} The prediction combines a cached DC prior with a
learned residual:
\begin{equation}
 \hat\theta_i=\operatorname{anchor}\left[
 \pi\tanh(\theta_i^{\rm DC}/\pi)+s_\theta\tanh(H_\theta(h_i)/s_\theta)
 \right].
\end{equation}
The magnitude head receives the start/setpoint magnitude explicitly:
\begin{equation}
 \hat v_i=\operatorname{clip}_{[v_i^{\min},v_i^{\max}]}
 H_v([h_i\Vert v_i^{\rm set}]).
\end{equation}
Finally, exact \(\pi\)-model formulas compute predicted line/transformer flows
from \((\hat v,\hat\theta,\hat P_g,\hat Q_g)\), and a separate feasibility head
uses 24 pooled stress features. The AC-PF adapter consumes only
\((\hat v,\hat\theta)\).

\paragraph{Gradient caveat.}
The DC prior is computed from \(\hat P_g\) through a detached CPU/SciPy path.
Under the default voltage-only PF loss, the generator head therefore receives no
gradient through the angle prior. Its pretrained behavior is retained, but PF
fine-tuning updates it only if a separate generator-head loss is enabled.

\section{AC-PF Training Loss and Metrics}

The voltage-state loss is
\begin{equation}
 \mathcal L_{\rm state}=\frac1N\sum_i
 \left[(\hat v_i-v_i^\star)^2+\wrap(\hat\theta_i-\theta_i^\star)^2\right].
\end{equation}
With \(\mathcal M_P=\mathcal V\setminus\mathcal V_{\rm REF}\) and
\(\mathcal M_Q=\mathcal V_{\rm PQ}\),
\begin{equation}
 \mathcal L_{\rm phys}=\frac{1}{|\mathcal M_P|+|\mathcal M_Q|}
 \left[\sum_{i\in\mathcal M_P}\rho(\Delta P_i)+
 \sum_{i\in\mathcal M_Q}\rho(\Delta Q_i)\right],
\end{equation}
where \(\rho(r)=\log\cosh(\operatorname{clip}(r,-30,30))\) by default. The
training objective is
\begin{equation}
 \boxed{\mathcal L=\mathcal L_{\rm state}+10^{-2}\mathcal L_{\rm phys}.}
\end{equation}
The script computes training physics through complex64 tensors but uses the
complex128 dataset/Ybus path for reported validation and test residual metrics.

\section{Interpretive Summary}
\begin{tabularx}{\linewidth}{>{\bfseries}p{0.25\linewidth}X}
\toprule
Property & GridSFM PF adaptation \\
\midrule
Graph & Seven node types: bus, proxy generator, load, shunt, AC branch, transformer branch, and cycle.\\
Topology encoding & Signed incidence, fundamental cycles, static Laplacian/DC features, and learned Hodge diffusion.\\
Neural core & Eight blocks of per-type linear attention, typed local GNN aggregation, and residual FFNs.\\
Voltage decoder & One terminal head; angle is DC prior plus learned residual, magnitude is bounded.\\
Physics inside forward & DC prior, exact branch-flow calculation, and terminal stress/feasibility features; no iterative AC residual feedback.\\
Physics outside forward & Exact-Ybus robust AC-PF loss and complex128 evaluation metrics.\\
Pretraining & Released GridSFM checkpoint may be loaded and fine-tuned.\\
\bottomrule
\end{tabularx}

\section*{Code provenance}
The formulation follows \path{gridsfm/model.py}, \path{gridsfm/blocks.py},
\path{gridsfm/cycle_basis.py}, \path{gridsfm/hodge_pe.py},
\path{gridsfm/pe_features.py}, \path{gridsfm/dc_prior.py}, and the AC-PF
adapter in \path{train_valid_test_gridsfm.py}, as present in August 2026.

\end{document}
